\begin{ouverture}[college1]{5}{Métaux : réactions et corrosion}{Partie 1 · Organisation et transformations de la matière}{Pourquoi certains métaux s’altèrent-ils ?}
\objectif{Exploiter une réaction acide-métal avec un modèle fourni.}
\objectif{Identifier gaz et ions à l’aide d’une banque.}
\objectif{Interpréter un protocole de corrosion et une protection.}
\prerequis{Équations fournies, ions et pH.}
\prerequis{Conservation de la masse.}
\end{ouverture}
\pcPage{Le cours}{Métaux : réactions et corrosion}
\phantomsection\label{ess-5}
\essSection{1}{Un acide peut réagir avec un métal}
Certains métaux réagissent avec une solution acide : du métal et des ions H+ sont consommés, des ions métalliques et du dihydrogène peuvent apparaître. Exemple fourni : $\ce{Zn + 2 H+ -> Zn^{2+} + H2}$. Les atomes et la charge totale sont conservés. Tous les métaux ne réagissent pas de la même manière avec un même acide ; on ne généralise pas un seul essai.\par
\pcFigure{figures/b05-reaction.pdf}{Modèle zinc-acide : les ions spectateurs ne figurent pas dans l’équation simplifiée.}
\essSection{2}{Interpréter les tests}
Un test du dihydrogène peut produire une petite détonation au contact d’une flamme, dans un dispositif pédagogique sécurisé par le professeur. Les ions formés peuvent être identifiés sur un autre prélèvement avec une banque de précipitation. La diminution des ions H+ peut s’accompagner d’une hausse du pH ; elle ne signifie pas forcément que la solution finale est neutre.\par
\clearpage
\essSection{3}{Comprendre la corrosion}
La corrosion est une transformation chimique d’un métal avec son environnement. La rouille du fer est un ensemble de composés ; sa formation habituelle nécessite notamment eau et dioxygène. Comparer des expériences en ne faisant varier qu’un facteur permet de rechercher les conditions. Une eau salée peut accélérer le phénomène ; il faut des mesures comparables pour quantifier cet effet.\par
\pcFigure{figures/b05-corrosion.pdf}{Expériences comparées : même métal, même durée ; résultats qualitatifs annoncés.}
\essSection{4}{Protéger et recycler}
Une peinture ou un revêtement isolant limite le contact du fer avec l’eau et le dioxygène tant que la barrière reste efficace. Certains métaux forment une couche superficielle protectrice. La protection dépend du matériau et du milieu. Recycler réduit certains prélèvements de ressources, mais demande tri et énergie. Ne pas confondre altération chimique, fusion du métal et simple fragmentation.\par
\begin{pcEncadre}{method}{Un exemple expliqué}Le fer réagit selon $\ce{Fe + 2 H+ -> Fe^{2+} + H2}$. On observe des bulles, puis le test du gaz indique H2 ; un test à la soude donne un précipité vert associé à Fe2+. Les deux résultats sont cohérents avec la réaction fournie. Les bulles seules n’auraient pas permis cette conclusion.\end{pcEncadre}
\pcPage{Expérimenter et retenir}{Métaux : réactions et corrosion}
\phantomsection\label{activite-5}
\begin{pcEncadre}{method}{Investigation en classe}Exploiter des photographies datées de clous identiques : air sec, eau au contact de l’air, eau salée au contact de l’air, eau désaérée protégée de l’air. Construire un tableau d’observations, choisir les comparaisons pertinentes et proposer une expérience pour tester un revêtement sans changer tous les paramètres.\end{pcEncadre}
\pcSousTitre{Ce que je dois savoir expliquer}\begin{questions}
\item Lire une équation acide-métal et repérer les espèces spectatrices.
\item Distinguer observation et identification du gaz.
\item Relier une protection aux conditions de corrosion.
\end{questions}
\begin{pcEncadre}{attention}{Une erreur à éviter}La rouille n’est pas « du fer sale ». C’est le résultat d’une transformation produisant d’autres espèces ; sa masse peut inclure des éléments venus de l’environnement.\end{pcEncadre}
\pcRepere{cor-activite-5}{Pistes pour l’investigation}\pcGrille
\clearpage\markright{Je démarre}
\rubrique{Je démarre}
\par\noindent\begin{minipage}{\linewidth}
\begin{exercice}[1]{Corrosion}
\begin{questions}
\item Le fer qui rouille subit-il une transformation physique ou chimique ?
\item Citer deux facteurs usuels nécessaires dans le modèle étudié.
\end{questions}
\end{exercice}
\end{minipage}
\par\noindent\begin{minipage}{\linewidth}
\begin{exercice}[1]{Un gaz}
\begin{questions}
\item Quel gaz est formé dans $\ce{Zn + 2 H+ -> Zn^{2+} + H2}$ ?
\item Quel ion métallique apparaît ?
\end{questions}
\end{exercice}
\end{minipage}
\par\noindent\begin{minipage}{\linewidth}
\begin{exercice}[1]{Les réactifs}
\begin{questions}
\item Dans $\ce{Fe + 2 H+ -> Fe^{2+} + H2}$, quelles espèces sont consommées ?
\item L’ion chlorure non écrit peut-il être spectateur ?
\end{questions}
\end{exercice}
\end{minipage}
\par\noindent\begin{minipage}{\linewidth}
\begin{exercice}[1]{Protection}
\begin{questions}
\item Quel rôle joue une peinture sur le fer ?
\item Pourquoi une rayure peut-elle réduire son efficacité ?
\end{questions}
\end{exercice}
\end{minipage}
\par\noindent\begin{minipage}{\linewidth}
\begin{exercice}[1]{Le bon test}
\begin{questions}
\item Une banque associe le précipité vert à Fe2+. Quel ion identifie-t-on ?
\item Pourquoi éviter de réaliser soi-même un test à la flamme à domicile ?
\end{questions}
\end{exercice}
\end{minipage}
\clearpage\markright{Je m’entraîne}
\rubrique{Je m’entraîne}
\par\noindent\begin{minipage}{\linewidth}
\begin{exercice}[2]{Conserver la charge}
\begin{questions}
\item Comparer la charge totale des deux côtés de $\ce{Zn + 2 H+ -> Zn^{2+} + H2}$.
\item Compter les atomes H de chaque côté.
\end{questions}
\end{exercice}
\end{minipage}
\par\noindent\begin{minipage}{\linewidth}
\begin{exercice}[2]{Le pH évolue}
\begin{questions}
\item Une réaction acide-fer fait passer le pH de 2 à 4. Que dire de l’acidité ?
\item Est-ce cohérent avec une consommation d’ions H+ ?
\end{questions}
\end{exercice}
\end{minipage}
\par\noindent\begin{minipage}{\linewidth}
\begin{exercice}[2]{Une pesée ouverte}
\begin{questions}
\item Un dispositif perd 0,06 g pendant une réaction produisant uniquement un gaz qui s’échappe. Quelle masse de gaz s’est échappée si aucune autre perte ne se produit ?
\item Pourquoi faut-il écarter les projections ?
\end{questions}
\end{exercice}
\end{minipage}
\par\noindent\begin{minipage}{\linewidth}
\begin{exercice}[2]{Choisir un couple}
\begin{questions}
\item A : air sec sans rouille ; B : eau + air avec rouille. Que suggère la comparaison ?
\item Quel essai supplémentaire testerait le rôle du dioxygène ?
\end{questions}
\end{exercice}
\end{minipage}
\par\noindent\begin{minipage}{\linewidth}
\begin{exercice}[2]{Ne pas tout changer}
\begin{questions}
\item Deux clous diffèrent par taille, métal, salinité et durée d’exposition. Peut-on isoler l’effet du sel ?
\item Décrire une meilleure comparaison.
\end{questions}
\end{exercice}
\end{minipage}
\clearpage\markright{Je raisonne}
\rubrique{Je raisonne}
\par\noindent\begin{minipage}{\linewidth}
\begin{exercice}[3]{Un bilan de masse}
\begin{questions}
\item Un métal de 5,6 g incorpore 2,4 g de matière venant du milieu. Quelle masse de produit si tout est conservé ?
\item Pourquoi cette augmentation ne viole-t-elle pas la conservation ?
\end{questions}
\end{exercice}
\end{minipage}
\par\noindent\begin{minipage}{\linewidth}
\begin{exercice}[3]{Une conclusion trop large}
\begin{questions}
\item Le cuivre ne produit pas de bulles dans une solution acide donnée pendant un essai. Peut-on conclure que « les acides ne réagissent jamais avec les métaux » ?
\item Formuler une conclusion limitée.
\end{questions}
\end{exercice}
\end{minipage}
\par\noindent\begin{minipage}{\linewidth}
\begin{exercice}[3]{Comparer des résultats}
\begin{questions}
\item Deux séries donnent des masses de fer corrodé de 0,20 g et 0,24 g après sept jours. L’écart seul prouve-t-il un effet du sel ?
\item Que faudrait-il pour rendre l’argument plus solide ?
\end{questions}
\end{exercice}
\end{minipage}
\par\noindent\begin{minipage}{\linewidth}
\begin{exercice}[3]{Choisir un modèle}
\begin{questions}
\item L’équation $\ce{Fe + 2 H+ -> Fe^{2+} + H2}$ décrit-elle à elle seule la rouille ordinaire dans l’air humide ?
\item Pourquoi distinguer les deux situations ?
\end{questions}
\end{exercice}
\end{minipage}
\par\noindent\begin{minipage}{\linewidth}
\begin{exercice}[3]{Réparer plutôt que jeter}
\begin{questions}
\item Proposer deux actions réduisant la consommation de métal pour un objet corrodé.
\item Pourquoi le recyclage ne dispense-t-il pas de toute consommation d’énergie ?
\end{questions}
\end{exercice}
\end{minipage}
\pcPage{Défi du chapitre}{Métaux : réactions et corrosion}
\phantomsection\label{defi-ex-5}
\pcDefiTitre{5}{Le portail près de la mer}
\begin{questions}
\item Un portail en fer rouille surtout aux rayures. Expliquer le rôle du revêtement.
\item Proposer une expérience isolant l’effet de la salinité.
\item Rédiger une recommandation justifiée et une limite de votre étude.
\end{questions}
\pcGrille
